% Part: methods
% Chapter: proofs
% Section: resources

\documentclass[../../../include/open-logic-section]{subfiles}

\begin{document}

\olfileid[id]{mth}{prf}{res}

\section{Sumber Lain}

Ada banyak buku mengenai cara menyusun bukti dalam matematika yang mungkin
berguna. Secara khusus, bacalah \emph{How to Read and do Proofs: An
Introduction to Mathematical Thought Processes} \citep{Solow2013} dan
\emph{How to Prove It: A Structured Approach} \citep{Velleman2019}.
\href{http://www.people.vcu.edu/~rhammack/BookOfProof/BookOfProof.pdf}{\emph{Book
of Proof}} \citep{Hammack2013} dan
\href{https://scholarworks.gvsu.edu/books/7/}{\emph{Mathematical
Reasoning}} \citep{Sandstrum2019} merupakan buku tentang bukti yang tersedia
secara bebas di internet. Para filsuf mungkin akan menganggap
\emph{More Precisely: The Math you need to do Philosophy}
\citep{Steinhart2018} sebagai pengantar yang baik untuk penalaran matematis.

Berbagai panduan yang lebih singkat tentang bukti juga tersedia di internet;
misalnya,
\href{https://math.berkeley.edu/~hutching/teach/proofs.pdf}{``Introduction
  to Mathematical Arguments''} \citep{Hutchings2003} dan
\href{https://eugeniacheng.com/wp-content/uploads/2017/02/cheng-proofguide.pdf}{``How to
  write proofs''} \citep{Cheng2004}.

\subsection{Video Motivasi}

Merasa tidak punya motivasi untuk mengerjakan tugas? Sedang merasa murung?
Video-video ini mungkin membantu!

\begin{itemize}
\item \url{https://www.youtube.com/watch?v=ZXsQAXx_ao0}
\item \url{https://www.youtube.com/watch?v=BQ4yd2W50No}
\item \url{https://www.youtube.com/watch?v=StTqXEQ2l-Y}
\end{itemize}

\end{document}
